\section{Translate the state transition into an array update}
The bit update has an outer byte loop and an inner loop of eight steps. At the outer
loop boundary, the state represents exactly the bytes already consumed. Check the
remaining length before loading another byte. With a zero length, return the incoming
state without dereferencing the pointer. Otherwise load one octet, XOR it into the
state, and reset the bit counter to eight.

In the inner loop, retain the old low bit before shifting. The branch decision must
refer to that retained bit; testing after the shift would select a different feedback
sequence. After eight steps, advance the full pointer by one and decrease the remaining
byte count. Return to the outer boundary. The counter cannot underflow on a valid
input because the zero test precedes every load and decrement.

The illustrated MIPS64 implementation keeps addresses in full registers and uses
\texttt{daddiu} to advance its byte walker. It uses \texttt{lbu} so the input
value is between zero and 255. An unmasked \texttt{lb} can introduce ones above the
low byte when bit seven is set, changing the state before its first shift.
The state uses logical doubleword shifts with an explicitly zero-extended feedback
constant. This keeps the 32-bit invariant visible in a 64-bit machine.

An implementation may use word instructions instead, but it must account for their
extension rules and return a zero-extended state. In particular, an arithmetic right
shift of a negative word fills high positions with ones. The unsigned CRC transition
requires zeroes. A correct low word with incorrect upper bits still violates this
project's return interface.

\subsection*{Exercise: Find a width error}
Use an input byte with bit seven set and an initial state with bit 31 set. Predict
which error a signed byte load can introduce and which error a signed word shift can
introduce. Choose observations before the feedback XOR that distinguish the two faults.
